\section{视觉算法功能设计}
在概念层面，视觉（导航）系统可拆分为七个相互解耦的算法功能模块：激光惯性里程计、高频位姿插值、定位健康监测、坐标变换与位姿转发、飞行指令生成、统一安全管线、飞行状态监测与安全夺权。各模块以独立 ROS 2 节点或独立函数单元的形式运行，通过话题解耦，任一模块失效均不会导致整体系统崩溃，从架构上保证了系统的鲁棒性与可调试性。

\subsection{激光惯性里程计（FAST-LIO）}
激光惯性里程计负责将机载 3D 激光雷达点云与 IMU 数据进行紧耦合融合，输出全局一致的里程计位姿。选用 FAST-LIO 方案：其基于迭代扩展卡尔曼滤波（iEKF），以 IMU 高频预测、雷达点云低频校正，兼具精度高、计算量小、对退化环境鲁棒的特点，适合机载算力受限与快速运动场景。里程计同时发布定位状态位标志（\lstinline{/localization_status}，bit0=位姿可用、bit1=导航可用），供控制与监控节点判断定位是否可信。

\subsection{高频位姿插值}
由于雷达里程计帧率较低（约$\qty{10}{\Hz}$），直接用于$\qty{50}{\Hz}$的控制回路会引入明显延迟。为此利用 IMU 数据对里程计进行插值，输出高频（$\qty{100}{\Hz}$量级）的平滑位姿，显著降低控制延迟与位姿抖动，是保证飞行控制品质的关键环节。

\subsection{定位健康监测}
为保证“永不坠机”，在位姿转发链路中内置健康监测逻辑：对里程计位置做数值有限性校验与坐标越界校验，对连续帧做跳变检测与帧超时检测。一旦发现异常即判定定位不可信并\textbf{锁存}该状态，停止向飞控转发视觉位姿，使飞控因位姿缺失自动退出视觉依赖，从而规避漂移位姿导致的失控风险。

\subsection{坐标变换与位姿转发}
将雷达里程计坐标系（机体 FLU、世界 ENU）变换为飞控期望的坐标系（MAVROS 内部约定 ENU、PX4 使用 NED），同步对速度、协方差矩阵做旋转映射，并通过 \lstinline{mavros} 转发至飞控。同时以 \lstinline{MAV_STATE_ACTIVE} / \lstinline{MAV_STATE_FLIGHT_TERMINATION} 状态字向飞控报告视觉定位进程的健康状况。

\subsection{飞行指令生成}
导航控制节点根据协议指令（\lstinline{std_msgs/UInt8MultiArray}）或遥控器拨杆选择飞行模式并生成控制指令，支持三种模式：
\begin{itemize}
    \item \textbf{方向模式（Direction）}：按前后左右/转向位标志输出恒定方向速度，并叠加直线保持、保高与保向修正；
    \item \textbf{航点模式（DirectSetpoint）}：对目标航点进行比例导航逼近，近场自动减速，到达后自动悬停并锁定航向；
    \item \textbf{反反无人机模式（AntiAntiDrone）}：在锚点悬停后于上下两个高度目标间周期切换，规避对空打击。
\end{itemize}
指令来源支持外部话题（协议）与遥控器预设航点（RC）双源，并可在运行中通过拨杆动态切换。

\subsection{统一安全管线}
无论指令来源于何处，所有输出指令均依次经过统一安全管线处理后再下发：有限性校验 $\rightarrow$ 速度限幅 $\rightarrow$ 加速度（斜坡）约束 $\rightarrow$ 地理围栏与最低高度约束 $\rightarrow$ 发布。该管线将机体运动能力与场地边界固化为软件约束，从根本上消除速度阶跃、超速与越界风险。

\subsection{飞行状态监测与安全夺权}
飞行状态监测节点将飞控连接、位姿可用、导航可用三级判定汇总为飞行状态（\lstinline{/drone/flight_status}）上报电控；遥控器摇杆监测节点在 OFFBOARD 模式下检测到操作手拨动摇杆时，主动请求飞控切换至 ALTCTL 模式，实现控制权一键交还，作为自动飞行失效时的最终人工兜底。

\section{自瞄系统设计}

空中机器人自瞄系统负责完成敌方机器人装甲板识别、空间定位、目标状态估计、运动预测、弹道解算以及云台与发射控制指令生成。针对空中平台高速运动、多种打击任务并行以及特殊目标运动状态复杂等特点，2026 赛季主要从\textbf{传感器时序、低延迟感知、目标运动建模、特殊目标处理、预测时延与多模式控制}等方面对自瞄系统进行完善，并进一步探索利用机体定位信息进行自身运动补偿。

\subsection{时序一致性与传感器数据组织}
空中机器人在实际飞行过程中会同时发生横滚、俯仰、偏航和平移运动，图像观测与姿态信息之间的时间误差会直接转化为目标空间定位误差。为此，传感器节点持续缓存带时间戳的 IMU 数据，并依据每帧图像的采样时刻在相邻姿态之间进行插值，使进入后续算法的图像和姿态具有统一时间基准。

在现有双相机方案基础上，本赛季进一步明确两套相机与不同攻击任务之间的数据边界。主相机负责读取外置 IMU，并通过共享话题提供统一姿态数据；第二相机使用同一姿态来源，但依据自身图像时间戳独立完成姿态匹配。两套视觉链分别维护相机标定、检测配置和预测参数，避免不同任务之间的状态与参数相互耦合。

\subsection{低延迟装甲板感知}
装甲板检测采用异步推理方式。检测节点完成图像预处理后将任务提交至独立推理模块，网络执行完成后再进行后处理和结果发布。系统限制同时处于推理状态的任务数量，当计算资源已经饱和时优先舍弃无法及时处理的新任务，而不是持续累积历史图像，从而保证预测器尽可能处理最新目标状态。

三维定位仍以装甲板四角点和相机标定为基础进行 PnP 解算，但目标观测在进入状态估计前进一步消除相机光心与云台旋转中心之间的固定偏移。这样可减少云台转动时由相机偏心运动引入的虚假目标速度，使运动状态更加接近目标相对于实际云台旋转中心的运动。

\subsection{目标运动模型与特殊目标处理}
系统继续按照不同目标类型选择对应运动模型，并在 2026 赛季进一步调整旋转目标的状态表示。普通小陀螺模型减少难以稳定估计的高阶加速度状态，主要使用目标中心位置、中心速度、旋转角和连续角速度描述整车运动，在保证预测能力的同时降低高阶状态噪声对滤波结果的影响。

针对前哨站，系统建立独立运动模型，显式维护目标中心、旋转角、角速度、旋转半径以及三块装甲板槽位之间的结构关系。前哨站处于静止状态时优先使用单装甲板模型；进入稳定旋转状态后再切换至专用旋转模型，避免静止目标受到不必要的高维状态估计误差影响。旋转角速度采用连续状态估计，并通过不同的进入、退出阈值形成迟滞，减少模型在临界状态附近频繁切换。

空中俯视前哨站时，顶部结构可能与腰部可击打装甲板同时进入视野。对此，系统分别建立正常腰部装甲板和顶部结构的三维几何假设，将模型角点重新投影至图像，并比较其与网络检测角点之间的重投影误差，以空间模型与实际图像的一致性完成顶部结构过滤。

\subsection{预测时延与多任务控制}
自瞄预测的目标并非当前时刻位置，而是控制指令真正执行并完成命中时目标所在的位置。因此系统按照图像处理、指令发送、电控响应、发弹和弹丸飞行等过程计算预测时间，并进一步将普通连续跟踪装甲板与前哨站等特殊目标的时延参数进行拆分，使不同模型能够根据自身运动形式和发射时机独立调整预测时刻。

当前系统同时运行普通装甲板自瞄、5 mm 自瞄、基地打击和能量机关等多条攻击链。各算法只向各自独立话题输出控制结果，最终由模式切换模块选择唯一有效控制源。模式切换时首先输出空闲指令并设置短暂保护时间；若当前有效算法在规定时间内没有继续更新控制结果，则自动回退至 \lstinline{IDLE}，避免过期控制量持续作用于云台和发射机构。

\subsection{高动态自身运动补偿探索}
飞机自身高速运动时，视觉测得的目标相对运动同时包含敌方真实运动与己方运动。针对这一问题，本赛季实现了一条基于定位里程计的 MAP 自运动补偿实验链：视觉继续提供敌方目标观测，定位模块提供飞机在统一地图坐标系中的位置、姿态和速度，系统分别预测目标与己方枪口在未来时刻的状态，再重新计算两者的相对瞄准关系。

\textbf{该链路目前仍处于实验与实车测试阶段，尚未作为正式比赛自瞄主链使用。}因此当前设计保留成熟视觉链与 MAP 实验链的可切换能力，在不影响正式自瞄系统稳定性的前提下继续验证高动态自身运动解耦的有效性。
